\documentclass[12pt,a4paper]{article}
\usepackage[scheme=plain]{ctex}
\usepackage[margin=2.5cm]{geometry}
\usepackage{indentfirst}
\usepackage{setspace}
\usepackage{fontspec}
\setCJKmainfont[BoldFont=SimHei,ItalicFont=KaiTi]{SimSun}
\newCJKfontfamily\kaifont{KaiTi}[AutoFakeBold]
\usepackage{newpxtext,newpxmath}

\usepackage{lastpage}





\pagestyle{plain}

\usepackage[nobottomtitles*]{titlesec}
\titleformat{\section}{\Large\bfseries}{}{0em}{ }
\titlespacing*{\section}{0pt}{3ex}{1.5ex}

\usepackage{xcolor}
\usepackage{needspace}

\setlength{\parindent}{2em}
\setlength{\parskip}{0.4em}

\doublespacing

\title{\bfseries Day 3\\[0.3em]村与城：如何接近远端的存在？}
\author{}
\date{}

\begin{document}
\maketitle

\vspace{-2em}

\vspace{1em}

\section{前哨的第三封信}
\vspace{-0.3em}{\color{gray}\small\kaifont\noindent K 在村庄继续推进他的工作。他试图靠近城堡——不止是从物理意义上靠近，更是从理解上靠近。他走到了村子的边缘，走进了一户陌生人的屋子，也走进了村公所。他看到了三种不同的东西，但他在信里没有告诉我们答案。他想问的是一个比答案更先的问题：如果你面对的是一个你看不到的东西，你该从哪里开始？}\vspace{0.5em}
\vspace{0.3em}

\noindent\rule{\textwidth}{0.4pt}

\vspace{0.3em}



同志们：



城堡在哪里？我已经从三个不同的方向试过了。结果互不相同。



先说最近的。昨天晚上我决定直接走过去。大路看起来是朝山上去的，走了很久，我确认了一件事：这条路永远不会到城堡。它在山脚下拐了弯，然后跟山平行，一直往前走，村子也不到头。不是“路走错了”那种感觉，是“路本身就是这么设计的”。



然后我进了一间屋子。屋子挺大的，有人洗澡，有人洗衣服，有个女人抱着婴儿坐在靠椅上，头也不抬，衣服泛着丝绸的光。她说是从城堡来的少女。其他人呢？不是我预想中的敌人也没有仇恨。老人给我推了一块木板让我从雪里爬出来，两个男人把我拽到门边送我出去，自始至终没有说一句“不可以”——但动作就是答案。在这间屋子里，城堡不在山顶上，它在每一个人看别人的方式里、在说话的次序里、在关门和开门的时机里。



最后我去了村公所。村长躺在床上，风湿病犯了。他讲了两件事：电话和公文。打到城堡的电话会让最下一级部门的全部电话一起响起来，但即便通了也不能当真——那是官员用来消遣的；公文在各部门之间来回传递，每个环节都有答复，但从没有裁断。



三样东西。空间上的隔绝，人际关系里的渗透，制度上的错位。它们指向的是同一个城堡吗？如果是，我该怎么把它们拼在一起？如果不是——那这个“城堡”到底是什么？



K

\newpage
\section{通往城堡的路}
\vspace{-0.3em}{\color{gray}\small\kaifont\noindent K 决定直接用双脚靠近城堡。大路看上去是朝山上去的，但走到山脚下时，路拐了弯，从此与山平行。他不是在这条路上迷了路，他是发现了这条路本身的目的。}\vspace{0.5em}
\vspace{0.3em}

\noindent\rule{\textwidth}{0.4pt}

\vspace{0.3em}



此刻，他在明澈的空气中看清了城堡的轮廓，且由于那层山上处处皆是薄薄的积雪把任何物体的形状都勾勒出来，形状益发明晰可辨了。山上的积雪看来比村里少得多。K此时在村里踏雪前进，并不比昨天在大路上省力。这里的积雪一直堆到与那些简陋小屋的窗户齐高，稍往上一点，低矮的屋顶上也沉甸甸地堆满了雪，而山上的那些建筑物，却都那么自由自在、轻松愉快地挺立着，至少从这里看上去的印象是如此吧。

站在这里，从远处看去，城堡的外观大体与K的预料相符。它既不是一座古老的骑士城堡，也不是一座新式的豪华建筑，而是一座宽阔的宫苑，其中两层楼房为数不多，倒是有许许多多鳞次栉比的低矮建筑；如果事先不知道这是一座城堡，那么可能会以为它不过是一个小镇。K只看见一座塔，它究竟是属于一所住宅还是属于一所教堂，就无法看清了。一群群乌鸦在围着它盘旋。

K心里什么别的都不想，两眼看准城堡继续朝前走。但是愈走近城堡，他就愈觉失望，那的确不过是一个相当寒酸的、看去全是一色普通村舍的小城镇，只有一个优点，就是或许所有的房子全是石结构的；但是墙上的灰泥早已掉光，砌墙的石块看来也开始剥落了。K蓦地想起了自己的家乡小镇；它同这座所谓的城堡相比几乎毫无逊色。如果K的目的只是观光，那么长途跋涉来到此地便太不值得，倒不如重访自己多年未归的故里更明智些呢。于是他在心中将家乡那座教堂的塔同眼前山上的塔作了一番比较。家乡那座塔从低到高、由粗而细笔直伸向天空，塔顶较宽，红瓦覆盖，诚然仍是一座不能超凡脱俗的建筑——我们凡人还能盖出什么别的样式来呢？——但它比大片低矮房屋有着更崇高的目标，比人们那晦暗的终日劳碌具有更明朗的意态。然而这里山上这座塔呢——这是此处唯一可见的塔，现在已能看清是一所住宅的塔，说不定就是城堡主楼的塔，它是一座单调的圆形建筑，部分蒙常春藤垂青加以覆盖，塔身有不少小窗子，此时在阳光下发出刺眼的反光——这使人觉得有些荒唐，塔顶类似阁楼，其雉堞瑟瑟缩缩、杂乱无章、残颓破败地戳向蓝天，就好像是一个害怕画错或是很不认真的孩子信手涂鸦胡乱画上去似的。面对这景象，人有一种感觉，仿佛这楼中有一个忧心忡忡的居民，按理本应老老实实地把自己关在楼中最偏僻的角落向隅而泣，现在居然冲破楼顶，探出身来向全世界亮相了。

K再次停下步来，似乎停下步子能增加他对眼前景象的判断能力。然而这时他受到了干扰。他站住的地方离村教堂（其实这不过是个小礼拜堂而已，只是像仓库似的扩建了一下，以便能容纳全村的教民）不远，教堂后面就是村子的小学。这学校是一座长长的低矮房子，奇怪的是它给人一种既是临时应急又非常古老的印象，楼前有一座围了铁栏杆、眼下是一片白雪覆盖着的花园。这时孩子们正好同老师一起从楼里走出来。一大群孩子把老师团团围住，每双眼睛都盯着他，从四面八方对他哇啦哇啦地唠叨个不住，K一点听不清他们那连珠炮似的话。这位老师是个年轻人，小个子，窄肩膀，不过样子并不可笑，身子直挺挺的，他从老远处一眼就看到了K，当然，除了他和他的一群学生之外，K是这方圆数里内目力所及唯一可见的人。K作为一个外乡人，首先打招呼问好，更何况是向这个在这里享有相当权威的小个子男人呢。“你好，老师！”他说。这话一出口，孩子们一下子就鸦雀无声，这突然的肃静作为他开口说话的先导，看样子使这位老师颇为得意。“您在观看城堡吗？”他问这问题时的态度，比K预料的要和蔼些，然而那语气似乎不赞同K这样做。“是的，”K说，“我不是本地人，是昨天晚上才到村里的。”——“您不喜欢这城堡吗？”教师很快又问。“什么？”K反问一句，有点震惊，然后用比较温和的语气重复那个问题：“我喜不喜欢这城堡吗？为什么您要猜想我不喜欢它呢？”——“没有一个外乡人喜欢它。”教师说。为了避免在这里说出一些不中听的话，K改变话题问：“您大概认识伯爵吧？”——“不认识。”教师说着便准备扭头走开。然而K并不罢休，他追问道：“什么？您不认识伯爵？”——“我怎么会认识他？”教师低声说，然后便用法语大声补充道：“请您考虑一下这里有这么些天真无知的孩子在旁边。”K听到这话便抓住时机问道：“老师，我可以拜访拜访您吗？我打算在这里待较长时间，可我现在就感到有点孤单了；我既不是农民中的一员，大概也不能算是城堡中的一员吧。”——“农民和城堡没有太大的区别。”教师说。“也许是这样，”K说，“可这丝毫改变不了我的处境。我可以拜访拜访您吗？”——“我住在天鹅巷肉铺。”虽然这话听起来只像是在说出住址，而不是发出邀请，K仍然说道：“好的，我一定来。”教师点点头，就同这时突然又大吵大嚷起来的孩子们一起继续走他的路了。不久，他们便消失在一条地势陡然下降的小巷里。

然而K现在却有些心神不宁，这次谈话使他心中颇感不快。自从到达此地以来，他现在第一次感到真正的疲倦了。到这里来的这一段很长的路，本来好像并不使他感觉吃力，这些天来，他是日复一日、从容不迫、一步一步地走过来的呵！——可是现在呢，过度紧张的后果到底还是显现出来了，无疑，来得真不是时候。他感到一种不可抗拒的强烈冲动，希望在这里尽量多结识一些人，但是每认识一个人都增添他一分倦意。〔1〕他想，在他今天这种竞技状态下，如果能咬牙坚持这次步行，至少走到城堡大门，那就已经是了不起的成绩了。

于是他又继续前行，但是路仍然很长。走着走着他发现，这条同时是村子主要街道的大路并不是通到城堡所在的山上去的，它只通到城堡近处，虽然眼看快到山脚下了，却像故意作弄人似的在那里拐了弯，然后，尽管沿它走下去并不会离城堡越来越远，却怎么也无法再接近它一步。K一直在期待着这条大路最终总会拐进城堡里去，也仅仅因为他抱有这一期望，他才不住地往前走；显然由于他感觉疲劳，才没有毅然决然离开它，另外，这个其长无比的村子也使他惊诧不已，它没有尽头，大路两边老是出现同样的小房子、冻了冰的窗户、厚厚的积雪，一个人影都不见——最后他终于还是下决心离开这条缠人磨人的大路，转进一条狭窄的小巷，这里雪更厚，把陷进雪里的脚拔出来艰巨异常，累得他满身大汗，突然间他站住了，再也走不动了。

细想起来，他倒也并不完全孤单，左右两旁不都是农舍吗？于是他捏了一个雪球，向一扇窗户掷去，门应声开了——这是他在村子里走了这大半天第一道打开的门，一个老农，身穿棕色皮夹克，头歪向一边，和蔼可亲地、颤巍巍地站在那里。“我可以到您屋里歇一会儿吗？”K说，“我太累了。”他完全没有听清老头说些什么，只是感激地接受了老者的好心：给他推过来一块木板，这木板当即把深陷雪中难以自拔的他救了出来，只跨过了三五步，他就进了老农屋里。
\par\vspace{0.3em}{\footnotesize\color{gray}（第一章）}
\newpage
\section{城堡来的少女}
\vspace{-0.3em}{\color{gray}\small\kaifont\noindent K 被老人从雪地里拉进了一间屋子。屋内光线昏暗，烟雾弥漫，有人洗澡，有人洗衣，有个女人靠在椅背上抱着婴儿。她说自己是"从城堡来的少女"。在这里，城堡不在山顶上，在每个人的行为方式里。}\vspace{0.5em}
\vspace{0.3em}

\noindent\rule{\textwidth}{0.4pt}

\vspace{0.3em}



屋子挺大，笼罩在一片朦胧的光线中。刚从外面进来的K一时什么也看不见。他趔趄几步，快撞到一个洗衣盆时，一只女人的手拉住了他。从屋子的一角传来好几个孩子的叫嚷声。从另一个角落则涌出滚滚烟雾，使本来就半明半暗的屋子变成一片漆黑。K如置身云雾中。“他喝醉了。”有人说。“您是谁？”另一个声音厉声喝问，然后大概是冲着老头问：“你为什么把他放进来？难道可以把大街小巷里东溜西窜的什么乱七八糟的人都放进来吗？”——“我是伯爵的土地测量员。”K说，试图用这话在那些直到现在仍唯闻其声不见其人的主人面前为自己辩护一下。“哦，原来是土地测量员。”这时响起的是一个女声，继之而来的是全然的寂静。“你们听说过我？”K问道。“当然。”同一个声音简短地答道。人家听说过他，这好像不是在抬举他吧。

浓烟终于稍稍散去了一些，K渐渐可以辨认四周了。看来今天是个大洗涤的日子。离门不远处有人在洗衣服，但雾气是来自屋子另一角，那里放着一个很大的木澡盆，K还从未见过这么大的澡盆——约有两张床那么大，两个男人泡在热气腾腾的水里洗澡。然而更加令人惊奇、可又弄不清究竟是什么东西使人惊奇的，是屋子右边那个角落。屋子后墙上有一个大缺口，通过这后墙上唯一的开口处，兴许是从院子里吧，射进来一道煞白的雪光，使一个坐在角落处一把高高的靠椅里、面带倦容的女人身上穿的衣服发出一种类乎丝绸的光亮。这女人怀里抱着一个婴儿，她四周有几个孩子在玩耍，显然都是农家孩子，但她似乎同他们并不属于一类人，当然，疾病和疲倦，也能使农民显得高雅些的。

“您请坐吧。”两个男人中的一个说，这人一脸络腮胡外加两髭向上翘起的小胡子，不断张着嘴喘气，他从水里伸出手，越过盆边——那副样子颇为滑稽——向一个大木箱指去，甩了K一脸热水。木箱上已经坐着一个人，正是那个把K放进来的老头，在那里发愣。对自己终于可以坐下来，K心里很是感激。现在没有人再理睬他了。在洗衣盆旁边的那个女人头发金黄，青春焕发，边洗边轻声哼着小曲，澡盆里的男人则不停地蹬腿翻身，孩子们想到他们跟前去，但一再被盆里溅出的大量洗澡水打退，K也未能幸免，靠椅里的女人泥塑木雕似的靠在椅背上，甚至连胸前的孩子她也不低头看上一眼，而只是视而不见地仰望空中。

K盯着她——这幅凝滞不动的、美丽、忧郁的雕像——大概看了好久，然后他一定是睡着了，因为当他听见有人大声叫他而猛然惊醒时，他的头是枕在他身旁的老者肩上的。这时两个男人早已洗完澡，穿好衣服站在K面前，澡盆里现在是孩子们在金发女人的看管下扑腾了。现在看来，两个男人中那个大嗓门的大胡子是地位低一些的。另一个的个头不比大胡子高，胡须也少得多，是个少言寡语、从容思考的人，他身材矮胖，长着一张宽脸，老是低着头。“土地测量员先生，”他说，“您不可以老待在这里，请原谅我的失礼。”——“我也不打算在这里待下去，”K说，“只是想稍稍休息一下。现在已经休息过了，我这就走。”——“您大概奇怪我们这里不太殷勤好客吧，”那人说，“我们这里可没有殷勤好客的习惯，我们是不需要客人的。”K小寐以后觉得精神稍好，注意力也比先前容易集中一些，听到这些直率的话很高兴。他行动不那么拘谨了，用手杖一会儿拄拄这，一会儿拄拄那，然后向靠椅里的女人走了过去。他也是这屋里个子最高的。

“不错，”K说，“你们要客人干什么呢。不过恐怕间或也会需要个把人的，比如我，土地测量员。”——“这个我不知道，”男人慢吞吞地说，“既然叫您来，那么大概是需要您，这可能是个例外，可我们呢，我们这些小老百姓，我们是按常规办事，这一点您不能怪我们。”——“哪里，哪里，”K说，“我对您只有感谢，对您和这里所有的人。”这时，出乎众人的意料，K突然嗖地一跳转过身去，站在那女人面前了。她用失神的蓝眼睛看着K，一块透明的真丝头巾一直盖到她前额正中，婴儿已在她怀里睡着了。“你是谁？”K问道。她拂袖——究竟这一轻蔑的表示是针对K的还是针对她自己的回答，弄不清楚——答道：“一个从城堡来的少女。”

这一切只是一瞬间的事。现在K发现自己已经一左一右夹在两个男人中间并被——好像根本没有别的交流思想的办法——一声不响地但却是用尽全力地拽到门边来。老头儿这时不知在乐什么，高兴得拍手。洗衣女人也在突然拼命吵嚷起来的孩子们旁边格格笑了。
\par\vspace{0.3em}{\footnotesize\color{gray}（第一章）}
\newpage
\section{电话与公文}
\vspace{-0.3em}{\color{gray}\small\kaifont\noindent K 来到村公所，村长卧病在床。他讲了两件事。打到城堡的电话会让最下一级部门的全部电话一起响起来，即便有人接听了也不能当真——那是"八音盒"。公文在各部门之间来回流转，每个环节都有回应，但最终没有裁断。如果说空间上的隔绝和人际关系里的渗透是你能感受到的痕迹，那制度上的错位呢？}\vspace{0.5em}
\vspace{0.3em}

\noindent\rule{\textwidth}{0.4pt}

\vspace{0.3em}



K对此地衙门的这种看法，首先在村长处得到了完全的证实。这是个和颜悦色、胡须刮得干干净净的胖子，他是在病中，在风湿病严重发作时躺在床上接见K的。“这么说您就是我们的土地测量员先生了。”他说，同时想坐起来表示欢迎，但怎么使劲也无用，于是带着歉意指指自己的腿又颓然躺回枕头上去。一个不苟言笑的女人给K拿来了一把椅子摆在床前，这间屋子的窗子本来就很小，又拉上了窗帘而更加昏暗，所以几乎只能看得见这女人那影影绰绰的轮廓。“您请坐，您请坐，土地测量员先生，”村长说，“请告诉我您有什么要求。”K念了克拉姆的信又加上几句自己的话。现在他再次体味到同衙门打交道时那种异常轻松的感觉。衙门简直就肩负着所有的重担，您可以把什么都推给它，自己落得个轻松自在。村长似乎从他那一方面同样感到了这一点，在床上很不舒服地翻了个身。最后他说：“您已经看到了，土地测量员先生，整个这件事情我是早就知道了的，我之所以还没有采取相应措施，第一是因为我生病，另外就是您很久没来，我已经在想您大概已经放弃这工作了。但现在呢，既然您亲临寒舍造访，我当然就必须把令人不快的全部实情倾囊相告。如您所说，您已被录用为土地测量员；但是遗憾得很，我们并不需要土地测量员。这里可以说一点让土地测量员做的工作也没有。我们这些小家小户的土地界线是划定了的，一应事项已全部有条不紊登记在案。产业易主的事极少发生，小的边界纠纷我们自己来解决。所以我们要土地测量员做什么？”虽然K原先完全没有细想过这一点，但他内心深处确信，自己是早就料到会得到类似说法了。〔7〕



正因为如此他立即说道：“我真万万没有想到事情是这样。这使我的全部打算都落空了。现在我只希望这是一个误会。”——“可惜没有误会，”村长说，“情况正像我说的那样。”——“可这也太不像话了！”K叫起来。“难道我千辛万苦长途跋涉到这里就为了让人再把我送回去吗！”——“这是另一码事，”村长说，“这件事我无权决定，但您所谓的那个误会是怎么产生的我倒可以给您解释解释。对一个像伯爵衙门这样的大衙门来说，有时难免会发生一个部门发出这一项指令、另一个部门发出另一项指令而互相不通气的事，虽然上一级的监督检查是极为严格的，但检查很自然地往往晚到一步，于是总会出点小差错，当然毫无例外都是在一些微乎其微的小事上出错，比如您这件事。在大事上我还没听说出过什么错，不过就是小事也够让人尴尬的了。至于说到您这件事，那么我可以坦率地把事情原原本本讲给您听，我用不着保守公务秘密——要保密我的官还太小，我是务农的，一辈子务农。很久以前，那时我当村长才几个月，上头来了一份公函，我记不清是哪个部门发出的了，那公函用高级长官特有的命令口气通知说，拟聘任一名土地测量员，指令村公所将有关该测量员工作必需的各项计划和全部文字材料立即备齐待命。这公函上指的人当然不可能是您，因为这是好多年前的事了，要不是现在我卧病在床，有的是时间去回想那些最最可笑的事情，那么我也是不会记起这事来的。”——“米齐，”他突然中断叙述对那个一直在屋里跑来跑去不知在忙乎什么的女人说，“请你去那边柜里看看，兴许能找到那份文件。”——“这项指令，”他接着对K解释，“是我刚当上村长那段时间下达的，那时我把所有的文件都保存起来了。”那女人当即打开柜子，K和村长在一旁看着她。柜子里各种文件纸张塞得满满当当的。柜门一开，两大捆像劈柴一样被捆成一大卷的文件便骨碌碌滚了出来，把那女人吓了一跳，一下子闪开了。“下面，可能在下面，”村长在床上指挥着，那女人则驯顺地张开双臂把大批大批的文件抱出来撂在地上，这样才好去拿最底下的那份文件。现在一卷卷文牍纸片已经堆满了半间屋子。“我们做了很多工作呵，”村长颔首说，“这些还只是一小部分呢。大宗的我都存放在库房里，而且，绝大部分文件都丢失了。谁有本事把全部文件都保存下来呀！库房里还有很多很多呢。”——“你找得着那份指令吗？”过一会儿他又对他女人说。“你先得找一份上面有‘土地测量员’字样，底下还画了蓝道道的文件。”——“这儿太暗了，”那女人说，“我去取一支蜡烛来，”说罢就迈步跨过地上的大堆文牍出屋去了。“我女人在这项繁重的公务中帮了我的大忙了，”村长说，“我还有别的重要事务，这项工作只能算是兼管一下。虽说我还有一个助手帮助我整理这些文字材料，这就是那位老师，可即使这样也仍然做不完，总是剩下许多堆积起来，还没整理的全部集中在那边那个柜子里，”他指指另一个柜子。“现在我这一病倒呢，那里简直就堆成山了，”说完这句他又疲倦地、然而却面带骄傲之色躺了下去。“我可不可以，”当那女人取来了蜡烛后又跪在柜子前继续找那份指令时，K说道，“帮着您太太一起找？”村长微笑着摇摇头：“刚才我说过，对您我没有什么公务秘密可保；可是，要让您本人去文件堆里找东西，那我确实又走得太远了。”此刻屋里一片安静，唯闻翻动纸张窸窸窣窣之声，在这种情况下村长甚至微微打起盹来。一阵轻轻的叩门声使K回身一看，当然又是两个助手。不过无论如何，两人现在总算有点长进：他们不是乒乒乓乓闯进屋来，而是站在外面通过微开的门缝轻声说：“我们在外面待着太冷了。”——“谁呀？”村长被惊醒了，问道。“没别人，是我的两个助手，”K说，“我不知道该让他们在哪儿等我，外面太冷，这儿嘛他们又让人讨厌。”——“我不在乎他们，”村长和气地说。“您让他们进来好了。其实我认识他们，我们是老相识了。”——“但是我讨厌他们，”K直言不讳，他把目光从两个助手移到村长身上，又从村长移回到助手身上，发现三个人的微笑一模一样，毫无区别。“既然你们两个到这儿来了，”他带着试探的口吻说，“那么就留下来帮着村长太太找一份文件吧，文件上写着‘土地测量员’，字底下画着蓝道道。”村长并无反对意见。就是说，不许K做的，两个助手倒可以做，两人也立即投身纸堆，甩开膀子大干起来，可是，与其说他们在找文件，倒不如说在纸堆里乱翻一气，而且，当一个人正读文件标题时，另一个人总是一把将文件从他手中抢走。那女人则跪在那现在已经空空如也的柜子前面，看样子根本没在继续找，至少现在蜡烛离她老远老远。



...



“您认识施瓦尔策吗？”K问。

“不认识，”村长说，“你大概认识吧，米齐？也不认识。不，我们不认识这个人。”

“这真是怪事，”K说，“他是城堡一位副主事的儿子。”

“亲爱的土地测量员先生，”村长说，“您怎么可以要求我认识所有城堡副主事的所有儿子？”

“很好，”K说，“那么您就只好相信我说的，承认他就是城堡一位副主事的儿子。同这个施瓦尔策，我在到这里的当天就有过一次不愉快的争论。后来他打电话给一个叫弗里茨的副主事询问，答复是我已经被聘任为土地测量员了。这您又怎么解释呢，村长先生？”

“很简单，”村长说，“这一点正好说明您从来还没有同我们的官府接触过嘛！所有这一切接触统统是表面文章，由于对情况完全无知，您竟将它们信以为真了。至于说到电话嘛，您瞧，我和官府在公务上的联系怎么说也是够多的了吧，可是我这里就没有电话。在酒店一类地方电话可能有相当大的好处，好比是八音盒那一类东西，然而更大的作用也就没有了。您在这里打过电话吗？打过？好，那您也许能明白我的意思。城堡里的电话显然非常灵；我听人说过，那里总是一个电话接着一个电话，这当然使工作的进度加快许多。那边连续不断的电话，我们这边电话机里听起来就是不停的嗡嗡声和歌唱声，这您一定也听到了。但我们这里的电话所能传达的唯一正确可信的东西，也就只是这种沙沙声和歌唱声罢了，别的信息全不是那么回事。我们同城堡之间并无确定的电话线路，没有总机接转我们打去的电话；从这里给城堡中某人挂电话时，那边最下一级的办事部门的电话就都一齐响起来，甚至于，这点我知道得很清楚，要不是城堡里几乎所有电话机上的响铃装置都关上了，那么所有的电话机就都会响起来的。偶尔有个别过于疲劳的官员觉着需要放松一下消遣消遣，特别是晚上或夜间，他就会把响铃装置打开；那时我们就能听到回话了，可那叫什么回话，不过是开开玩笑而已。这也是完全合情合理的。谁有权利为了想解决自己一点小小的私人困难而用电话铃声去干扰那些极为重要的、忙得人晕头转向的紧张工作？我也不明白为什么连一个外乡人也会以为，要是他比如说打电话给索尔蒂尼，那么接电话答复他的就一定是索尔蒂尼了。说多半会是另一个毫不相干的部门的一个小小记录员，这可能性不是反倒更大些吗？反过来说，在极少有的时候当然也会出现打电话给小记录员时反倒是索尔蒂尼本人来接电话的情况。那时当然是趁着还没听到他开腔说第一句话就赶紧跑开为妙。”
\par\vspace{0.3em}{\footnotesize\color{gray}（第五章）}
\newpage
\section{城市偏向（urban bias）：划一条区分“城”与“村”的线}
\vspace{-0.3em}{\color{gray}\small\kaifont\noindent 在发展中国家，城市与乡村之间那条线能不能划、怎么划，是理解一切城乡问题的前提。在城堡的世界里，我们面临一个类似的问题：城堡与村庄之间也同样找不到明确的边界。或许，我们不能因为没有边界就放弃分析。}\vspace{0.5em}
\vspace{0.3em}

\noindent\rule{\textwidth}{0.4pt}

\vspace{0.3em}



本章旨在界定“城市偏向”概念，并就其作为贫困与发展共存现象的罪魁祸首提出一个初步论证（后续章节将予以充实），同时附带说明它也是发展本身诸多低效与不足的成因。至此，我们已界定了“偏向”一词，并考察了两个我们期望据以证明大多数穷国存在“城市偏向”的规范：效率规范和公平规范。就前者而言，我们必须证明，如果将更大份额的资源投向农村，产出（按照第52页时间加权的含义）会提高。要依据公平规范来证明城市偏向，我们在上一节中已经论证过，大概只需证明农村人口比城市人口致富更慢（这在一定程度上是政策所致），以及农村人口起初就比城市人口贫穷很多，便已足够。后续章节将确立农村相对被剥夺的事实，考察农村资源以超出效率或公平合理限度的方式被转移至城市的情况，并检视导致这类转移的压力。



首先，我们必须仔细审视“城乡”二分法本身。可能有四种批评意见。（1）“这条线无法划定。城市在不知不觉中过渡为农村，中间经由城市份地、花园郊区与乡镇企业作为中介，界线模糊不清。”（2）“不止一条线。两个部门太少了。具体而言，存在一个第三类‘城市’部门，主要提供各类服务，在决定何种资源配置最优——或判断实际配置是否存在偏向时——必须加以考虑。”（3）“这条线应该画在别处：要画在首都与全国其他地方之间，而非画在城市与农村地区之间。”（4）“这条线画错了地图。资源配置与偏向存在于农业与工业之间，或消费品与投资品之间，或穷人与富人之间，而非农村与城市之间。”



**城乡界线**



在穷国，通常可以在城市与农村之间画出一条相当清晰的界线。少数例外不难想到：廉价且高度发达的交通系统使得兼职务农与进城通勤成为可能的地区，如斯里兰卡的湿润区；拥有可观小规模园艺产业的半发达第三世界城邦，如新加坡；市场与行政中心密布、近乎连续农村聚落的地区，如印度南部的喀拉拉邦。然而，一般而言，穷国既有城镇景观亦有乡村景观，二者之间的断裂是鲜明的。穷人很少有钱或有足够的热量能量进行长途通勤，住在农村（或住在城市）的人往往也就在当地工作。



第三世界农村地区的三种主要生活兼劳作模式是游牧、自营农场或庄园农场务农，以及定居村庄。在前两种情形中，从事露天农业的正是同一批人，他们也采用农村的生活方式。这使得农村部门的边界非常清晰。在第三种情形下，村庄与城镇之间的界线偶尔会模糊。一座城镇通常具有若干特征（紧凑性、人口密度、规模、非农业主导性）和诸多功能（市场、行政、社会、教育、交通），人口在三四千到七八千之间的某些地方满足部分标准但不满足另一些。不过，这个问题无论在什么意义上都是边缘性的。



穷国中绝大多数村庄社群的规模是数百人，主要是农业人口或为农业服务的工匠。大多数城镇人口居住在一万几千人以上的地方，极少从事农业或直接服务于农业的工作。因此，尽管各国人口普查对“城市”的定义看似杂乱多样，但对于绝大多数地方和居民而言，城乡二分法依然是清晰且离散的。本书所诊断的偏向源于人口一万到两万及以上规模的较大城镇，并使这类城镇从中受益。



\vspace{0.3em}

\noindent\rule{\textwidth}{0.4pt}

\vspace{0.3em}
\par\vspace{0.3em}{\footnotesize\color{gray}（Lipton, M. (1977). \textit{Why poor people stay poor: Urban bias in world development}. Temple Smith. 由 DeepSeek V4 Pro 翻译。）}
\end{document}